% Generated draft: requires mathematical and source review.
\documentclass[11pt]{article}
\usepackage[margin=25mm]{geometry}
\usepackage{fontspec}
\setmainfont{Latin Modern Roman}
\usepackage{amsmath,amssymb,mathtools,unicode-math}
\setmathfont{Latin Modern Math}
\usepackage{enumitem,needspace,xurl,hyperref}
\setlength{\parindent}{0pt}
\setlength{\parskip}{5pt}
\setlength{\emergencystretch}{3em}
\newcommand{\N}{\mathbb{N}}
\newcommand{\Z}{\mathbb{Z}}
\newcommand{\Q}{\mathbb{Q}}
\newcommand{\R}{\mathbb{R}}
\newcommand{\C}{\mathbb{C}}
\newcommand{\F}{\mathbb{F}}
\newcommand{\A}{\mathbb{A}}
\newcommand{\PP}{\mathbb{P}}
\newcommand{\E}{\mathbb{E}}
\newcommand{\eps}{\varepsilon}
\newcommand{\dd}{\,\mathrm{d}}
\begin{document}
\section*{UnsolvedMath ID: 30004046}
\textbf{Source problem code:} catalogue entry
\subsection{Definitions and mathematical statement}
The supplied record poses three questions about localized list-coloring bounds, but does not identify the underlying theorems or their hypotheses. Consequently, it does not determine a fully quantified conjecture.

To describe the terminology conditionally, let \(G=(V,E)\) be a finite simple undirected graph. This graph class is an explanatory convention here; it is not specified in the supplied statement. For \(v\in V\), write
\[
N_G(v)=\{u\in V:\{u,v\}\in E\},\qquad d_G(v)=|N_G(v)|,\qquad N_G[v]=N_G(v)\cup\{v\}.
\]
A clique is a set of pairwise adjacent vertices, and \(\omega(H)\) is the maximum size of a clique in \(H\), with \(\omega(\varnothing)=0\). Two possible meanings of the undefined local clique number are
\[
\omega_G^{\mathrm{closed}}(v)=\omega(G[N_G[v]])
\quad\text{and}\quad
\omega_G^{\mathrm{open}}(v)=\omega(G[N_G(v)]).
\]
The first is the largest size of a clique containing \(v\); these two quantities differ by one. The record does not specify which convention its notation \(\omega(v)\) uses.

A list assignment \(L\) assigns each vertex \(v\) a set \(L(v)\) of permitted colors. A proper \(L\)-coloring is a function \(c:V\to\bigcup_{v\in V}L(v)\) such that \(c(v)\in L(v)\) for every vertex and \(c(u)\ne c(v)\) whenever \(\{u,v\}\in E\). For a prescribed function \(b:V\to\N\), a list-size bound would ordinarily assert that every list assignment satisfying \(|L(v)|\ge b(v)\) for all \(v\in V\) admits such a coloring. No particular function \(b\), graph hypotheses, or parameter quantifiers are supplied here.

\textbf{Source questions, with unresolved qualifications.}
\begin{enumerate}
\item For the intended local list versions of graph-coloring theorems, what conclusions hold when the local clique number \(\omega(v)\) is close to \(d_G(v)\)? Neither the meaning of ``close'' nor a proposed conclusion is given.
\item Is the intended restriction that the range of vertex degrees be at most exponential necessary? The record supplies neither a definition of this degree range nor the parameter in which the bound is exponential.
\item Can the intended list-coloring bounds be localized simultaneously? The record does not identify the bounds to combine, the quantities to localize, or the resulting list-size inequality.
\end{enumerate}
These questions retain the supplied scope, but their missing hypotheses and targets prevent replacement by a precise theorem or conjecture without further source evidence.
\subsection{Short English statement}
The problem asks whether local list-coloring results extend to vertices whose local clique number is close to their degree, whether an exponential restriction on degree variation is necessary, and whether several bounds can be localized simultaneously. The supplied record leaves the underlying theorems, bounds, and parameter conventions unspecified, so a fully quantified conjecture cannot yet be stated.
\subsection{Sources}
\begin{itemize}
\item Supplied local UnsolvedMath registry record, permanent ID \(30004046\), title \emph{Simultaneously Localized List-Coloring Bounds}, exact source problem code \texttt{catalogue entry}. The record assigns proposed year 2019 and supplies no proposer attribution.
\item Supplied source citation: \emph{Graph Theory (2020)}, Oberwolfach Reports, DOI 10.4171/owr/2019/1, pp.~5--63. \url{https://doi.org/10.4171/owr/2019/1}. The background describes these reports as records of mathematical workshops and their open-problem discussions. The supplied literature list calls this source \emph{Oberwolfach Report 1/2019}; the differing year descriptions are retained without resolution.
\item Supplied literature reference: Bonamy, Kelly, Nelson, and Postle, local coloring. \url{https://doi.org/10.1016/j.jctb.2022.06.002}.
\item Supplied literature reference: \emph{A local epsilon version of Reed conjecture}. \url{https://arxiv.org/abs/1911.02672}.
\item Supplied literature assessment, checked on \textbf{2026-08-22}: the record labels the problem open and states that several local list-coloring extensions are proved, while the questions concerning near-clique degree, degree range, and simultaneous localization remain open. This is supplied dated curation derived from the OWR record, not an independently verified literature assessment; time-sensitive status and cited claims require independent verification.
\item The registry labels the statement \texttt{corrected\_verified} and says that mathematical content was preserved while notation, typography, duplicate clauses, or omitted definitions were repaired for publication. This is a supplied editorial label, not mathematical verification of this draft or resolution of its missing hypotheses.
\end{itemize}
\end{document}
